% ==============================================================
%  Naawbi LMS --- CS3009 Software Engineering (Spring 2026)
%  Deliverable 1  |  Team Ibwaan
% ==============================================================
\documentclass[12pt, a4paper]{article}

% ── Packages ──────────────────────────────────────────────────
\usepackage[T1]{fontenc}
\usepackage[utf8]{inputenc}
\usepackage{lmodern}
\usepackage[top=2.5cm, bottom=2.5cm, left=2.5cm, right=2.5cm, headheight=28pt]{geometry}
\usepackage{microtype}
\usepackage{setspace}
\usepackage{parskip}
\usepackage{titlesec}
\usepackage{titling}
\usepackage{fancyhdr}
\usepackage{graphicx}
\usepackage{booktabs}
\usepackage{longtable}
\usepackage{array}
\usepackage{tabularx}
\usepackage{multicol}
\usepackage{enumitem}
\usepackage{mdframed}
\usepackage{xcolor}
\usepackage{hyperref}
\usepackage{marvosym}
% Provide a simple GitHub placeholder icon
\newcommand{\faGithub}{\raisebox{-0.1ex}{\Keyboard}}
\usepackage{soul}
\usepackage{float}
\usepackage{caption}
\usepackage[most]{tcolorbox}

% ── Colours ───────────────────────────────────────────────────
\definecolor{primary}{HTML}{1A237E}   % deep indigo
\definecolor{accent}{HTML}{283593}    % slightly lighter indigo
\definecolor{light}{HTML}{E8EAF6}     % very light indigo tint
\definecolor{rule}{HTML}{5C6BC0}      % mid indigo for rules
\definecolor{muted}{HTML}{546E7A}     % blue-grey for captions
\definecolor{rowalt}{HTML}{F3F4FA}    % alternating table row
\definecolor{boxbg}{HTML}{EEF0FB}     % info box background
\definecolor{success}{HTML}{1B5E20}   % deep green

% ── Hyperref setup ────────────────────────────────────────────
\hypersetup{
    colorlinks   = true,
    linkcolor    = primary,
    urlcolor     = accent,
    citecolor    = primary,
    pdftitle     = {Naawbi LMS — Deliverable 1},
    pdfauthor    = {Team Ibwaan},
    pdfsubject   = {CS3009 Software Engineering Spring 2026},
    pdfkeywords  = {LMS, Software Engineering, Naawbi, Ibwaan},
}

% ── Section formatting ────────────────────────────────────────
\titleformat{\section}
  {\Large\bfseries\color{primary}}
  {\thesection}{1em}{}
  [\color{rule}\titlerule]

\titleformat{\subsection}
  {\large\bfseries\color{accent}}
  {\thesubsection}{1em}{}

\titleformat{\subsubsection}
  {\normalsize\bfseries\color{accent}}
  {\thesubsubsection}{1em}{}

\titlespacing*{\section}{0pt}{1.6ex plus .5ex minus .2ex}{0.8ex}
\titlespacing*{\subsection}{0pt}{1.0ex plus .3ex minus .1ex}{0.4ex}

% ── Header / Footer ───────────────────────────────────────────
\pagestyle{fancy}
\fancyhf{}
\renewcommand{\headrulewidth}{0.4pt}
\renewcommand{\headrule}{\hbox to\headwidth{\color{rule}\leaders\hrule height \headrulewidth\hfill}}
\fancyhead[L]{\small\color{muted}\textbf{Naawbi LMS} — Deliverable 1}
\fancyhead[R]{\small\color{muted}CS3009 Spring 2026}
\fancyfoot[C]{\small\color{muted}\thepage}

% ── List styling ──────────────────────────────────────────────
\setlist[itemize]{leftmargin=1.5em, itemsep=2pt, topsep=4pt}
\setlist[enumerate]{leftmargin=1.5em, itemsep=2pt, topsep=4pt}

% ── Graceful image fallback ───────────────────────────────────
% Usage: \imgOrPlaceholder[width=...]{path}{caption text}
\newcommand{\imgOrPlaceholder}[3][width=0.85\textwidth]{%
  \IfFileExists{#2}{%
    \includegraphics[#1]{#2}%
  }{%
    \begin{center}%
    \fboxsep=8pt\fbox{\parbox{0.7\textwidth}{\centering\color{muted}%
      \textit{[Placeholder --- image not yet available]}\\[4pt]%
      \footnotesize #3}}%
    \end{center}%
  }%
}

% ── Custom environments ───────────────────────────────────────
\tcbuselibrary{skins, breakable}

\newtcolorbox{infobox}[1][]{
  enhanced, breakable,
  colback=boxbg, colframe=rule,
  arc=4pt, boxrule=0.8pt,
  leftrule=4pt,
  title=#1,
  fonttitle=\bfseries\color{primary},
  #1
}

\newtcolorbox{userstory}[1][]{
  enhanced, breakable,
  colback=white, colframe=rule!40,
  arc=3pt, boxrule=0.5pt,
  top=4pt, bottom=4pt, left=8pt, right=8pt,
  #1
}

% ──────────────────────────────────────────────────────────────
\begin{document}
\onehalfspacing

% ═══════════════════════════════════════════════════════════════
%  TITLE PAGE
% ═══════════════════════════════════════════════════════════════
\begin{titlepage}
  \centering

  \vspace*{2.5cm}

  % University & course
  {\small\scshape\color{muted} National University of Computer and Emerging Sciences}\\[0.4em]
  {\small\color{muted} CS3009 --- Software Engineering \quad|\quad Spring 2026}

  \vspace{1.5cm}
  {\color{rule}\hrule height 1.5pt}
  \vspace{1.5cm}

  % Project title
  {\fontsize{56}{62}\selectfont\bfseries\color{primary} Naawbi}\\[0.6em]
  {\Large\color{muted} Learning Management System}

  \vspace{1.2cm}
  {\color{rule!40}\rule{5cm}{0.6pt}}
  \vspace{1.2cm}

  {\LARGE\bfseries\color{primary} Deliverable 1}\\[0.4em]
  {\normalsize\color{muted} Requirements \& Sprint 1 Documentation}

  \vspace{2cm}

  % Team metadata
  \begin{tabular}{r@{\quad}l}
    {\color{muted}\small Team Name}       & {\bfseries Ibwaan} \\[5pt]
    {\color{muted}\small Team Lead}       & {\bfseries Ibraheem Farooq} \\[5pt]
    {\color{muted}\small Members}         & Ibraheem Farooq,\; Wajih-Ur-Raza Asif,\; Ali Aan \\[5pt]
    {\color{muted}\small Submission Date} & 28 February 2026 \\[5pt]
    {\color{muted}\small Repository}      & \href{https://github.com/aliaankhowaja/Naawbi}{github.com/aliaankhowaja/Naawbi} \\
  \end{tabular}

  \vfill
  {\color{rule!40}\hrule height 0.5pt}
  \vspace{0.6cm}
  {\footnotesize\color{muted}
    \href{https://github.com/ibbi1020}{ibbi1020}\enspace\textbullet\enspace
    \href{https://github.com/wajih-rathore}{wajih-rathore}\enspace\textbullet\enspace
    \href{https://github.com/aliaankhowaja}{aliaankhowaja}
  }
  \vspace{0.8cm}
\end{titlepage}
\nopagecolor

% ── Table of Contents ─────────────────────────────────────────
\tableofcontents
\newpage

% ═══════════════════════════════════════════════════════════════
%  SECTION 1 — INTRODUCTION
% ═══════════════════════════════════════════════════════════════
\section{Introduction}

\subsection{Context}
This project is a university Software Engineering team project to build a Learning
Management System (LMS). The system will implement standard SE practices
(requirements, design, implementation, testing, and documentation) and deliver a
working product with an MVP scope plus selected extensions.

\subsection{Problem Statement}
Current course management for small-to-mid scale classes often suffers from
fragmented tools for announcements, content, submissions, grading, and learner
progress visibility. Students lack a clear, data-driven view of performance and
deadlines. Instructors need a single workflow to publish content, collect
submissions, grade with rubrics, and monitor class progress.

\subsection{Goals}
\begin{itemize}
  \item Deliver a usable LMS that supports the end-to-end workflow for a course:
        \textbf{content} $\rightarrow$ \textbf{assignments} $\rightarrow$
        \textbf{submissions} $\rightarrow$ \textbf{grading}.
\end{itemize}

% ═══════════════════════════════════════════════════════════════
%  SECTION 2 — PROJECT VISION
% ═══════════════════════════════════════════════════════════════
\section{Project Vision}

Build a clean, reliable, course-centric LMS that instructors can use to manage
learning materials and evaluations, and that students can use to track progress
and performance. The key differentiators are:

\begin{itemize}
  \item \textbf{Grading} — rubric-driven grading for assignments.
  \item \textbf{Analytics Dashboard} — performance, progress, and deadline
        dashboards for both students and faculty.
\end{itemize}

% ═══════════════════════════════════════════════════════════════
%  PROJECT TEAM INFORMATION
% ═══════════════════════════════════════════════════════════════
\section{Project Team Information}

\noindent
\begin{tabular}{r@{\quad}l}
  \textbf{Team Name}    & Ibwaan \\
  \textbf{Project Name} & Naawbi \\
  \textbf{Team Lead}    & Ibraheem Farooq \\
\end{tabular}

% ── Role Table ────────────────────────────────────────────────
\subsection{Team Roles}

\begin{center}
\begin{tabularx}{0.9\linewidth}{>{\bfseries}l l X}
  \toprule
  Member & Role(s) & GitHub \\
  \midrule
  Ibraheem Farooq     & PM, Scrum Master          & \href{https://github.com/ibbi1020}{github.com/ibbi1020} \\[4pt]
  Wajih-Ur-Raza Asif  & Requirements Architect, Tester & \href{https://github.com/wajih-rathore}{github.com/wajih-rathore} \\[4pt]
  Ali Aan             & Developer, UI Designer    & \href{https://github.com/aliaankhowaja}{github.com/aliaankhowaja} \\
  \bottomrule
\end{tabularx}
\end{center}

% ── Member Bios ───────────────────────────────────────────────
\subsection{Member Biographies}

\medskip

% ---- Ibraheem Farooq ----------------------------------------
\noindent
\begin{minipage}[t]{0.30\linewidth}
  \vspace{0pt}
  \imgOrPlaceholder[width=\linewidth]{assets/ibraheem_photo.png}{Photo}
\end{minipage}%
\hspace{1em}%
\begin{minipage}[t]{0.73\linewidth}
  \vspace{0pt}
  {\bfseries\color{primary} Ibraheem Farooq}\\[-2pt]
  {\small\color{muted} PM / Scrum Master}\\[4pt]
  A Computer Science major based in Islamabad who transitioned from designing
  posters and brand identities to UX and product design. His technical background
  provides a unique perspective on bridging digital experiences that balance user
  needs with technical viability. Passionate about solving real problems through
  thoughtful design and clean, intuitive interfaces.
\end{minipage}

\medskip\noindent\color{rule!25}\hrule\color{black}\medskip

% ---- Wajih-Ur-Raza Asif -------------------------------------
\noindent
\begin{minipage}[t]{0.30\linewidth}
  \vspace{0pt}
  \imgOrPlaceholder[width=\linewidth]{assets/wajih_photo.jpeg}{Photo}
\end{minipage}%
\hspace{1em}%
\begin{minipage}[t]{0.73\linewidth}
  \vspace{0pt}
  {\bfseries\color{primary} Wajih-Ur-Raza Asif}\\[-2pt]
  {\small\color{muted} Requirements Architect / Tester}\\[4pt]
  A software engineer focused on building practical systems that solve real
  operational problems. His background spans end-to-end projects across different
  technologies, with a strong interest in automation, AI-driven workflows, and
  scalable system design. Currently expanding his understanding of enterprise
  technology and systems that make processes more efficient and reliable.
\end{minipage}

\medskip\noindent\color{rule!25}\hrule\color{black}\medskip

% ---- Ali Aan ------------------------------------------------
\noindent
\begin{minipage}[t]{0.30\linewidth}
  \vspace{0pt}
  \imgOrPlaceholder[width=\linewidth]{assets/ali_photo.png}{Photo}
\end{minipage}%
\hspace{1em}%
\begin{minipage}[t]{0.73\linewidth}
  \vspace{0pt}
  {\bfseries\color{primary} Ali Aan}\\[-2pt]
  {\small\color{muted} Developer / UI Designer}\\[4pt]
  An AI Engineer driven by a simple goal: building applications that work at
  scale. His expertise lies at the intersection of high-performance backend
  architecture and agentic systems. He specialises in designing and deploying
  scalable infrastructures that power intelligent applications. Beyond the code,
  he is a robotics enthusiast passionate about bringing new ideas to life.
\end{minipage}

% ── Team Agreement ────────────────────────────────────────────
\subsection{Team Agreement}

\subsubsection*{Communication}
\begin{itemize}
  \item \textbf{WhatsApp} --- primary channel for daily coordination; respond within 1 hour.
  \item \textbf{In-person} --- preferred when on campus.
  \item \textbf{Google Meet} --- for remote meetings.
  \item Urgent issues: tag the team and respond within 20 minutes.
  \item If unavailable: post an async update in WhatsApp within 24 hours.
\end{itemize}

\subsubsection*{Meetings}
\begin{itemize}
  \item Mandatory: Sprint Planning, Sprint Review, Sprint Retrospective.
  \item Regular check-ins: 4$\times$ per week (in-person or Google Meet).
  \item Absence: notify team 24 hours in advance and share an async update.
  \item Every meeting ends with: decisions + action items + owners.
  \item Notes: rotating assigned person.
\end{itemize}

\subsubsection*{Meeting Preparation}
\begin{itemize}
  \item \textbf{Sprint Planning:} review backlog, propose task breakdowns.
  \item \textbf{Review:} prepare demo flow + screenshots.
  \item \textbf{Retrospective:} each member prepares (1)~what went well,
        (2)~what did not, (3)~what to change.
\end{itemize}

\subsubsection*{Version Control}
\begin{itemize}
  \item Small, atomic commits; new branch per ticket.
  \item PR required before merging; do not commit secrets, credentials, or large binaries.
  \item Commit message format: \texttt{feat:}, \texttt{fix:}, \texttt{docs:}, \texttt{chore:}
\end{itemize}

\subsubsection*{Division of Work}
\begin{tabularx}{\linewidth}{>{}l X}
  \toprule
  \textbf{Owner} & \textbf{Responsibility} \\
  \midrule
  Wajih     & Requirement ownership: user stories, structured specs, and consistency checks. \\[3pt]
  Ali       & UI ownership and feature implementation. \\[3pt]
  Ibraheem  & Coordination, sprint tracking, packaging, and final submission. \\
  \bottomrule
\end{tabularx}

\medskip
\noindent\textbf{Contingency Plan.}
If a member becomes unavailable or falls behind, the team will re-scope to the
minimum required deliverable content, redistribute work immediately, and
prioritise required artefacts. The PM/SM will update the board, assign new owners,
and schedule an urgent meeting to unblock progress.

\subsubsection*{Submission Protocol}
\begin{itemize}
  \item \textbf{Submission owner:} Wajih.
  \item \textbf{Internal review:} Ibraheem and Ali.
  \item \textbf{Internal cutoff:} finalise 24 hours before the deadline.
\end{itemize}

% ═══════════════════════════════════════════════════════════════
%  SECTION 3 — INTENDED USE OF THE SYSTEM
% ═══════════════════════════════════════════════════════════════
\section{Intended Use of the System}

\subsection{Stakeholders}

\subsubsection*{Students}
\begin{itemize}
  \item View enrolled courses, announcements, and deadlines.
  \item Submit assignments.
  \item View grades and instructor comments.
  \item Use the analytics dashboard to understand:
    \begin{itemize}
      \item current standing and grade breakdown;
      \item missing/late work;
      \item upcoming deadlines and workload.
    \end{itemize}
\end{itemize}

\subsubsection*{Instructors}
\begin{itemize}
  \item Create and manage courses.
  \item Create assignments with due dates and submission rules.
  \item Grade submissions and provide feedback.
  \item View course-level insights (basic analytics):
    \begin{itemize}
      \item submission rates;
      \item grade distributions;
      \item at-risk indicators.
    \end{itemize}
\end{itemize}

% ═══════════════════════════════════════════════════════════════
%  SECTION 4 — FEATURES / OVERALL FUNCTIONALITY
% ═══════════════════════════════════════════════════════════════
\section{Features and Overall Functionality}

\subsection{What the System Will Do}

\subsubsection{Authentication and Role-Based Access Control (RBAC)}
\begin{itemize}
  \item User self-registration with role selection (Student or Instructor).
  \item Login and logout.
  \item Roles: Instructor, Student (optional: Teaching Assistant).
  \item Access control enforced across all features.
\end{itemize}

\subsubsection{Course Enrollment and Management}
\begin{itemize}
  \item Course creation (instructor).
  \item Course catalogue/listing.
  \item Enrollment workflow: students enroll using a randomly generated,
        unique code per class.
  \item Course home: stream of announcements and assignments, similar to
        Google Classroom.
\end{itemize}

\subsubsection{Announcements and Course Material}
\begin{itemize}
  \item Instructors post announcements visible to enrolled students.
  \item Announcement feed per course.
  \item File-based course material upload and distribution.
\end{itemize}

\subsubsection{Assignments and Submissions}
Assignments are a sub-category of announcements with the added functionality
of file upload and resubmission rules (similar to Google Classroom).
\begin{itemize}
  \item Instructors create assignments with title, description, due date/time,
        and late policy.
  \item Students submit via file upload with configurable resubmission rules
        (e.g., allowed until due date).
  \item Threaded comments within assignments.
  \item Submission tracking: timestamped submissions with status
        (\emph{submitted / late / missing}).
\end{itemize}

\subsubsection{Grading, Feedback, and Gradebook}
\begin{itemize}
  \item Instructor grading interface: view submissions, assign scores, add
        comments.
  \item Gradebook:
    \begin{itemize}
      \item Instructor view — all students and graded items.
      \item Student view — personal grades and feedback.
    \end{itemize}
\end{itemize}

\subsubsection{To-Do List and Deadlines \texorpdfstring{\normalfont\small(Key Extension)}{}}
A simple list within the course catalogue section providing student-facing
insights:
\begin{itemize}
  \item \textbf{Completion/progress indicators:} submitted vs.\ missing items.
  \item \textbf{Deadline awareness:} upcoming due dates and late/missed items.
\end{itemize}

\subsection{How the System Helps Users}
\begin{itemize}
  \item Consolidates course operations (content, announcements, assessments,
        grading) into one unified workflow.
  \item Reduces instructor workload through rubric reuse.
  \item Provides traceability for evaluation: all actions are recorded and
        auditable (who submitted what and when; who graded what and when).
\end{itemize}

% ═══════════════════════════════════════════════════════════════
%  SECTION 5 — USER STORIES
% ═══════════════════════════════════════════════════════════════
\section{User Stories}

\subsection{Authentication and Role-Based Access Control (RBAC)}
\begin{enumerate}[label=\textbf{US\,\arabic*}]
  \item As a \emph{User} (Student/Instructor), I want to log in securely so
        that I can access my account and authorised system features.
  \item As a \emph{User}, I want to log out of my account so that my session
        is secure when I am not using the system.
  \item As a \emph{User}, I want to select my role (Student or Instructor)
        during registration so that I can use the platform appropriately.
  \item As a \emph{User}, I want role-based access control to be enforced
        across all features so that individuals can only perform actions
        permitted by their chosen role.
\end{enumerate}

\subsection{Course Enrollment and Management}
\begin{enumerate}[label=\textbf{US\,\arabic*}, start=5]
  \item As an \emph{Instructor}, I want to create a new course so that I have
        a dedicated space to manage learning materials and track student
        progress.
  \item As a \emph{User}, I want to view a course catalogue/listing so that I
        can see the available courses in the system.
  \item As an \emph{Instructor}, I want the system to generate a unique random
        code for my course so that I can share it securely with my target
        students.
  \item As a \emph{Student}, I want to enroll in a course using a unique class
        code so that I can self-register for the classes I am supposed to take.
  \item As a \emph{Student}, I want to view a course home page that displays
        announcements and upcoming items so that I am immediately aware of
        current course activities.
\end{enumerate}

\subsection{Announcements and Course Material}
\begin{enumerate}[label=\textbf{US\,\arabic*}, start=10]
  \item As an \emph{Instructor}, I want to post announcements so that I can
        broadcast important updates or instructions to all enrolled students.
  \item As a \emph{Student}, I want to see an announcement feed for my course
        so that I never miss critical information from my instructor.
  \item As an \emph{Instructor}, I want to upload and publish course materials
        so that my students have access to the necessary learning resources.
  \item As a \emph{Student}, I want to access and download course materials so
        that I can study and complete my assignments.
\end{enumerate}

\subsection{Assignments and Submissions}
\begin{enumerate}[label=\textbf{US\,\arabic*}, start=14]
  \item As an \emph{Instructor}, I want to create assignments with a title,
        description, and due date/time so that expectations are clear to
        students.
  \item As an \emph{Instructor}, I want to define late policies and
        resubmission rules so that I can enforce grading standards.
  \item As a \emph{Student}, I want to submit my assignments via file upload
        so that I can turn in my completed work electronically.
  \item As an \emph{Instructor}, I want the system to timestamp all
        submissions so that I can accurately determine their status (submitted,
        late, missing).
  \item As a \emph{Student}, I want to see the status of my submissions (e.g.,
        submitted on time, late, missing) so that I know my work has been
        received.
\end{enumerate}

\subsection{Grading, Feedback, and Gradebook}
\begin{enumerate}[label=\textbf{US\,\arabic*}, start=19]
  \item As an \emph{Instructor}, I want to view a list of student submissions
        for an assignment so that I can review them efficiently.
  \item As an \emph{Instructor}, I want to assign scores and add comments to
        student submissions so that I can provide them with constructive
        feedback.
  \item As an \emph{Instructor}, I want to view a comprehensive gradebook of
        all students and their graded items so that I can monitor the overall
        performance of the class.
  \item As a \emph{Student}, I want to view my personal gradebook with scores
        and instructor feedback so that I can track my academic standing and
        progress.
\end{enumerate}

\subsection{To-Do List and Deadlines \texorpdfstring{\normalfont\small(Key Extension)}{}}
\begin{enumerate}[label=\textbf{US\,\arabic*}, start=23]
  \item As a \emph{Student}, I want to view a to-do list within my course
        catalogue section so that I have complete visibility over my progress
        (submitted vs.\ missing items).
  \item As a \emph{Student}, I want the system to highlight upcoming due dates
        and late/missed items so that I can effectively manage my deadlines and
        prioritise my workload.
\end{enumerate}

% ═══════════════════════════════════════════════════════════════
%  SECTION 6 — STRUCTURED SPECIFICATIONS
% ═══════════════════════════════════════════════════════════════
\section{Structured Specifications for All User Stories}
\label{sec:specs}

\subsection{Authentication and Role-Based Access Control (RBAC)}

\begin{userstory}
  \textbf{US\,1 \& US\,2 — Secure Login and Logout}\\[4pt]
  \textit{Roles:} Student, Instructor\\[2pt]
  \textit{Description:} Users must be able to authenticate themselves to
  access the system safely, and end their session when finished.\\[4pt]
  \textit{Acceptance Criteria:}
  \begin{itemize}
    \item System presents a login form requiring email and password.
    \item System authenticates credentials against the database.
    \item Successful login redirects to the user's respective dashboard.
    \item System provides a visible logout button that terminates the session
          and redirects to the login page.
  \end{itemize}
\end{userstory}

\begin{userstory}
  \textbf{US\,3 \& US\,4 — Registration and Role Enforcement}\\[4pt]
  \textit{Roles:} Student, Instructor\\[2pt]
  \textit{Description:} Users can register an account, select their role,
  and the system consistently enforces access restrictions based on this
  role.\\[4pt]
  \textit{Acceptance Criteria:}
  \begin{itemize}
    \item Registration page includes a required field to select ``Student'' or
          ``Instructor'' role.
    \item Students attempting to access instructor-only routes (e.g., course
          creation, grading) receive an ``Access Denied'' error.
    \item Instructors cannot access student-only views if not enrolled as a
          student.
  \end{itemize}
\end{userstory}

\subsection{Course Enrollment and Management}

\begin{userstory}
  \textbf{US\,5 — Course Creation}\\[4pt]
  \textit{Roles:} Instructor\\[2pt]
  \textit{Description:} Instructors can create courses to manage their
  classes.\\[4pt]
  \textit{Acceptance Criteria:}
  \begin{itemize}
    \item Instructor can access a ``Create Course'' form.
    \item Form requires Course Name and optional description/syllabus.
    \item Upon creation, the course is added to the instructor's dashboard.
  \end{itemize}
\end{userstory}

\begin{userstory}
  \textbf{US\,6 \& US\,9 — Course Catalogue and Home Page}\\[4pt]
  \textit{Roles:} Student, Instructor\\[2pt]
  \textit{Description:} Users can view a list of their courses and access a
  centralised course home page featuring announcements and upcoming
  items.\\[4pt]
  \textit{Acceptance Criteria:}
  \begin{itemize}
    \item Dashboard displays all active courses the user is enrolled in or
          teaching.
    \item Clicking a course opens the Course Home with a feed of recent
          announcements, latest assignments, and immediate deadlines (similar
          to Google Classroom's Stream).
  \end{itemize}
\end{userstory}

\begin{userstory}
  \textbf{US\,7 \& US\,8 — Unique Class Code Enrollment}\\[4pt]
  \textit{Roles:} Student, Instructor\\[2pt]
  \textit{Description:} The system generates a unique join code for courses,
  which students use to self-enroll.\\[4pt]
  \textit{Acceptance Criteria:}
  \begin{itemize}
    \item Upon course creation, a unique alphanumeric code is automatically
          generated.
    \item Instructor's course settings page clearly displays the join code.
    \item Student dashboard has a ``Join Course'' button accepting a valid
          class code as input.
    \item Entering a valid code adds the student to the course roster
          immediately.
  \end{itemize}
\end{userstory}

\subsection{Announcements and Course Material}

\begin{userstory}
  \textbf{US\,10 \& US\,11 — Course Announcements}\\[4pt]
  \textit{Roles:} Student, Instructor\\[2pt]
  \textit{Description:} Instructors publish text/media updates that populate
  a centralised feed for students.\\[4pt]
  \textit{Acceptance Criteria:}
  \begin{itemize}
    \item Instructors have an input box on the Course Home to post an
          announcement.
    \item Announcements support basic text formatting.
    \item Upon posting, the announcement is instantly visible on the Course
          Home feed for all students in the course.
  \end{itemize}
\end{userstory}

\begin{userstory}
  \textbf{US\,12 \& US\,13 — Course Material Distribution}\\[4pt]
  \textit{Roles:} Student, Instructor\\[2pt]
  \textit{Description:} Instructors upload relevant class materials for
  student retrieval.\\[4pt]
  \textit{Acceptance Criteria:}
  \begin{itemize}
    \item Instructors can attach files (PDFs, docs, links) to posts or a
          dedicated ``Materials'' section.
    \item Students can view or download the attached files without errors.
  \end{itemize}
\end{userstory}

\subsection{Assignments and Submissions}

\begin{userstory}
  \textbf{US\,14 \& US\,15 — Assignment Creation and Policies}\\[4pt]
  \textit{Roles:} Instructor\\[2pt]
  \textit{Description:} Instructors define tasks with explicit instructions,
  deadlines, and submission rules.\\[4pt]
  \textit{Acceptance Criteria:}
  \begin{itemize}
    \item Instructor can create an assignment specifying Title, Description,
          Total Points, Due Date, and Due Time.
    \item Instructor can toggle late submission permissions (e.g., ``Allow
          late turn-ins'' vs.\ ``Hard deadline'').
    \item Assignment appears in the course feed and on students' To-Do lists
          once published.
  \end{itemize}
\end{userstory}

\begin{userstory}
  \textbf{US\,16, US\,17 \& US\,18 — Assignment Submissions and Tracking}\\[4pt]
  \textit{Roles:} Student, Instructor\\[2pt]
  \textit{Description:} Students submit files, and the system records the
  exact turn-in timestamp, comment history, and completion status.\\[4pt]
  \textit{Acceptance Criteria:}
  \begin{itemize}
    \item Students have a dedicated ``Submit'' area within each assignment to
          upload files.
    \item System records the exact server time of the upload.
    \item Submissions are automatically tagged as:\
      \begin{itemize}
        \item \textbf{Submitted} — if before deadline;
        \item \textbf{Late} — if after deadline (provided late policy allows);
        \item \textbf{Missing} — if deadline passed with no upload.
      \end{itemize}
    \item Students can unsubmit/resubmit if the deadline has not passed and
          instructor settings allow it.
    \item Both students and instructors can leave threaded or inline comments
          on an assignment.
  \end{itemize}
\end{userstory}

\subsection{Grading, Feedback, and Gradebook}

\begin{userstory}
  \textbf{US\,19 \& US\,20 — Instructor Grading Interface}\\[4pt]
  \textit{Roles:} Instructor\\[2pt]
  \textit{Description:} Instructors evaluate student submissions, assign
  numeric scores, and leave qualitative comments.\\[4pt]
  \textit{Acceptance Criteria:}
  \begin{itemize}
    \item Instructors can open a list of all enrolled students for a specific
          assignment.
    \item Instructors can quickly view or download associated submission files.
    \item Instructors can input a numeric grade (up to the maximum defined
          points) and a text comment.
    \item Grades and feedback are saved and immediately reflected for the
          respective student.
  \end{itemize}
\end{userstory}

\begin{userstory}
  \textbf{US\,21 \& US\,22 — Comprehensive Gradebook Visibility}\\[4pt]
  \textit{Roles:} Student, Instructor\\[2pt]
  \textit{Description:} Instructors see a macro-view of class performance,
  while students see their individual progress.\\[4pt]
  \textit{Acceptance Criteria:}
  \begin{itemize}
    \item Instructor sees a grid/table where rows are students and columns are
          assignments, with total grades calculated.
    \item Students access a personal ``Grades'' tab showing all their graded
          assignments, earned points, and instructor feedback in one list.
  \end{itemize}
\end{userstory}

\subsection{To-Do List and Deadlines \texorpdfstring{\normalfont\small(Key Extension)}{}}

\begin{userstory}
  \textbf{US\,23 \& US\,24 — Student Progress Tracking}\\[4pt]
  \textit{Roles:} Student\\[2pt]
  \textit{Description:} Students receive automated insights into their
  workload, deadlines, and missing work.\\[4pt]
  \textit{Acceptance Criteria:}
  \begin{itemize}
    \item Student dashboard features a ``To-Do'' section outlining all pending
          assignments across all enrolled courses.
    \item The ``To-Do'' list is sorted chronologically by approaching
          deadlines.
    \item Passed deadlines without submissions are prominently labelled as
          ``Missing''.
    \item Items are automatically removed from the active To-Do list once
          submitted by the student.
  \end{itemize}
\end{userstory}

% ═══════════════════════════════════════════════════════════════
%  SECTION 7 — SCRUM BOARD
% ═══════════════════════════════════════════════════════════════
\section{Scrum Board}

\subsection{Sprint 1 Module}

\begin{center}
\begin{tabularx}{0.8\linewidth}{lX}
  \toprule
  \textbf{Module Name} & Course Management \\
  \textbf{Sprint Goal} & Implement core course management features to get the ball rolling. \\
  \bottomrule
\end{tabularx}
\end{center}

\subsection{Sprint 1 Selected User Stories}

\begin{figure}[H]
  \centering
  \imgOrPlaceholder[width=0.85\textwidth]{../assets/sprint1_user_stories.png}{Sprint 1 User Stories}
  \caption{Sprint 1 Selected User Stories (Trello screenshot)}
\end{figure}

\subsection{Four Structured Specifications (Sprint 1)}

The following structured specifications from Section~\ref{sec:specs} are
selected for Sprint 1:
\begin{enumerate}
  \item \textbf{US\,1 \& US\,2:} Secure Login and Logout
  \item \textbf{US\,3 \& US\,4:} Registration and Role Enforcement
  \item \textbf{US\,5:} Course Creation
  \item \textbf{US\,7 \& US\,8:} Unique Class Code Enrollment
\end{enumerate}

\subsection{Scrum Board Snapshots (Trello)}

\begin{figure}[H]
  \centering
  \imgOrPlaceholder[width=0.9\textwidth]{../assets/trello_snapshot_start.png}{Trello Start}
  \caption{Trello Board — Start of Sprint 1}
\end{figure}

\begin{figure}[H]
  \centering
  \imgOrPlaceholder[width=0.9\textwidth]{../assets/trello_snapshot_mid.png}{Trello Mid}
  \caption{Trello Board — Mid-Sprint}
\end{figure}

\begin{figure}[H]
  \centering
  \imgOrPlaceholder[width=0.9\textwidth]{../assets/trello_snapshot_end.png}{Trello End}
  \caption{Trello Board — End of Sprint 1}
\end{figure}

\subsection{Sprint 1 Product Burn-Down Chart}

\begin{figure}[H]
  \centering
  \imgOrPlaceholder[width=0.85\textwidth]{../assets/sprint1_burndown_chart.png}{Burndown Chart}
  \caption{Sprint 1 Burn-Down Chart}
\end{figure}

% ═══════════════════════════════════════════════════════════════
%  SECTION 8 — NON-FUNCTIONAL REQUIREMENTS
% ═══════════════════════════════════════════════════════════════
\section{Non-Functional Requirements (NFR) Specification}
\label{sec:nfr}

\begin{center}
\begin{tabularx}{\linewidth}{>{\bfseries\color{primary}}l X}
  \toprule
  NFR Category & Requirement \\
  \midrule
  Performance
    & Primary pages (course home, module list, gradebook summary) load within
      2–3 seconds under normal class sizes (\texttildelow 200 students, a few
      courses) on typical lab hardware. \\[6pt]
  Security
    & RBAC enforced so students cannot access other students' private
      grades/submissions. User inputs are validated and sanitised to prevent
      common web vulnerabilities. \\[6pt]
  Reliability
    & Submissions are not lost once acknowledged. Errors are handled
      gracefully with actionable error messages. \\[6pt]
  Usability
    & Consistent UI across modules (common navigation and layout). Clear
      status indicators for submissions (\emph{submitted / late / missing}). \\[6pt]
  Maintainability
    & Layered architecture (controller / service / repository). Unit tests for
      core business logic; at least basic integration tests for critical
      flows. \\[6pt]
  Scalability
    & Supports multiple courses and concurrent users typical of a
      department-scale deployment (not enterprise scale). \\[6pt]
  Auditability
    & Important events (logins, submissions, grading updates, enrollment
      changes) are logged for traceability. \\
  \bottomrule
\end{tabularx}
\end{center}

% ═══════════════════════════════════════════════════════════════
%  SECTION 9 — ITERATION 1 IMPLEMENTATION SCREENSHOTS
% ═══════════════════════════════════════════════════════════════
\section{Iteration 1 Implementation Screenshots}

\begin{figure}[H]
  \centering
  \imgOrPlaceholder[width=0.9\textwidth]{../assets/iteration1_1.png}{Iteration 1 Screenshot 1}
  \caption{Iteration 1 — Screenshot 1}
\end{figure}

\begin{figure}[H]
  \centering
  \imgOrPlaceholder[width=0.9\textwidth]{../assets/iteration1_2.png}{Iteration 1 Screenshot 2}
  \caption{Iteration 1 — Screenshot 2}
\end{figure}

\begin{figure}[H]
  \centering
  \imgOrPlaceholder[width=0.9\textwidth]{../assets/iteration1_3.png}{Iteration 1 Screenshot 3}
  \caption{Iteration 1 — Screenshot 3}
\end{figure}

% ═══════════════════════════════════════════════════════════════
%  SECTION 10 — WORK DIVISION
% ═══════════════════════════════════════════════════════════════
\section{Work Division}
\label{sec:workdiv}

\subsubsection*{Ibraheem Farooq \normalfont\small\color{muted}(PM / Scrum Master)}
\begin{itemize}
  \item Sprint planning, sprint tracking, and team coordination.
  \item Maintains delivery checklist for required artefacts and final submission packaging.
  \item Oversees branch/PR merge discipline and release readiness.
\end{itemize}

\subsubsection*{Wajih-Ur-Raza Asif \normalfont\small\color{muted}(Requirements Architect)}
\begin{itemize}
  \item Requirements ownership: user stories, structured specifications, and consistency checks.
  \item Verifies Sprint 1 scope alignment with requirements and acceptance criteria.
  \item Supports validation of deliverable completeness.
\end{itemize}

\subsubsection*{Ali Aan \normalfont\small\color{muted}(Developer / UI Designer)}
\begin{itemize}
  \item Development ownership for Sprint 1 scope (feature implementation).
  \item UI design ownership and implementation of user-facing screens.
  \item Produces Iteration 1 screenshots and demo-ready flows.
\end{itemize}

% ──────────────────────────────────────────────────────────────
\end{document}